\subsection{Access Atomicity and Alignment}
Unaligned (:term:base.terms.ordinary:) normal-memory accesses are architecturally permitted. A completed unaligned load may observe bytes from
multiple (:term:base.terms.memory_event|plural:) in coherence order, and a completed unaligned store may enter coherence as separately visible portions. Each
visible store portion is coherent and multi-copy atomic.

An unaligned store is nevertheless fault-atomic for synchronous faults. If any byte of the complete store range
would raise a synchronous fault, no byte of the store becomes architecturally visible.

Atomic read-modify-write instructions require natural byte alignment. A misaligned atomic operand raises
(:ref:base.events.EXCEPTION.ATOMIC_ALIGNMENT_FAULT:) before any memory update or architectural result
commits; it must not complete as a non-atomic update.

For every scalar size, a naturally aligned normal-memory load or store is tear-free. Unaligned accesses retain the
successful-access and synchronous-fault rules above.

Natural alignment equals the operand width in bytes: B is one byte and may use any byte address; W, L, and Q are
two, four, and eight bytes and require addresses divisible by two, four, and eight respectively.

Stores to memory that is later executed as instructions are (:term:base.terms.ordinary:) data stores until an explicit synchronization
sequence makes the modified bytes visible to instruction fetch.
